% Lösungen Einheit 1 von 5 – Monotonie und erste Ableitung (zusammenbau v0.1)
\einheitenkopf{Einheit 1 von 5 · Monotonie und erste Ableitung}

\erg{9}{a) steigt – die Tangente in $P$ geht nach rechts oben. \quad b) $f$ steigt auf diesem Intervall – für $x \ge 4$ ist $f'(x) = 2x - 6 \ge 2$, also positiv. \quad c) $f'(x) = 3x^2 - 6x - 9 = 0$ für $x_1 = -1$ und $x_2 = 3$; $f$ steigt auf $]-\infty; -1]$, fällt auf $[-1; 3]$ und steigt auf $[3; \infty[$. \quad d) $f'(x) = -3x^2 + 6x + 9 = 0$ für $x_1 = -1$ und $x_2 = 3$; $f'(-2) = -15 < 0$: $f$ fällt auf $]-\infty; -1]$; $f'(0) = 9 > 0$: $f$ steigt auf $[-1; 3]$; $f'(4) = -15 < 0$: $f$ fällt auf $[3; \infty[$. \quad e) $f$ steigt auf $]-\infty; -2]$, fällt auf $[-2; 1]$ und steigt auf $[1; \infty[$. \quad f) $f'(x) = x^2 \ge 0$ für alle $x$; $f'$ ist nur an der einzelnen Stelle $x = 0$ null, also ist der Graph monoton steigend. \quad g) $f'(x) = 3x^2 + 2 \ge 2 > 0$ für alle $x$; $f'$ wird nie null, es gibt keine waagerechte Tangente und damit keinen Extrempunkt. \quad h) $d(x) = \frac{1}{4}x^3 - 3x + 5$; $d'(x) = \frac{3}{4}x^2 - 3 = 0$ für $x = 2$ (im Intervall); $d$ fällt auf $[0; 2]$ und steigt auf $[2; 4]$; $d(2) = 1$, $d(0) = 5$, $d(4) = 9$; Wertebereich $[1; 9]$ – das Maximum liegt am rechten Rand.}
\erg{10}{a) Fehler: Die Monotonie wurde am Vorzeichen von $f''$ statt von $f'$ abgelesen. Richtig: $f'(x) = 3x^2 - 12x = 3x(x - 4) = 0$ für $x_1 = 0$, $x_2 = 4$; $f$ steigt auf $]-\infty; 0]$, fällt auf $[0; 4]$ und steigt auf $[4; \infty[$. \quad b) Ein Quadrat ist nie negativ; dazu kommt eine positive Zahl, also ist $f'(x)$ immer mindestens $1$ und damit nie null. \quad c) $A'(t) = -0{,}3t^2 + 3t = 0$ für $t = 10$, $A$ nimmt auf $[0; 10]$ zu, also $10$ Tage; $B'(t) = -0{,}6t^2 + 4{,}8t = 0$ für $t = 8$, $B$ nimmt auf $[0; 8]$ zu, also $8$ Tage. Die Aussage ist wahr.}
